\section*{第 \ref{ch:g8:electric-current-effects} 章——电流及其效应}

\begin{solution}{exo:g8:electric-current-effects:1}
\omterm{prop:g8:electric-current-effects:three}{热效应}（烤面包机、电水壶、灯丝）；\omterm{prop:g8:electric-current-effects:three}{磁效应}（
\omterm{def:g8:electric-current-effects:electromagnet}{电磁铁}、\omterm{ex:g6:circuits-and-safety:motor}{电动机}、门铃）；\omterm{prop:g8:electric-current-effects:three}{化学效应}（电镀、
电解水、给\omterm{def:g3:first-electric-circuit:battery}{电池}充电）。
\end{solution}

\begin{solution}{exo:g8:electric-current-effects:2}
\omterm{prop:g8:electric-current-effects:three}{热效应}；\omterm{prop:g8:electric-current-effects:three}{磁效应}；\omterm{prop:g8:electric-current-effects:three}{化学效应}；\omterm{prop:g8:electric-current-effects:three}{热效应}；\omterm{prop:g8:electric-current-effects:three}{磁效应}。
\end{solution}

\begin{solution}{exo:g8:electric-current-effects:3}
电流能在导线周围产生磁性——这是电与磁之间的第一座
桥梁。电流方向反转，
小磁针的偏转方向也随之反转。
\end{solution}

\begin{solution}{exo:g8:electric-current-effects:4}
\omterm{def:g8:electric-current-effects:electromagnet}{电磁铁}可以随时断电、通电，还能随意改变\omterm{def:g4:magnets-and-compass:poles}{磁极}方向，
其磁性强弱也可以调节（增加匝数、增大
电流）。永久\omterm{def:g1:magnets:magnet}{磁铁}则以上这些都做不到。
\end{solution}

\begin{solution}{exo:g8:electric-current-effects:5}
若裸露的线圈匝直接相互接触，电流就会在绕组间\omterm{def:g7:short-circuits-safety:device}{短路}，
而不再流经整个线圈——因此必须绝缘，迫使电流沿着
长长的磁化路径流动。另外，线圈本身对电流的阻碍很小：
若直接接在\omterm{def:g3:first-electric-circuit:battery}{电池}两端，几乎就是\omterm{def:g7:short-circuits-safety:device}{短路}，线圈和\omterm{def:g3:first-electric-circuit:battery}{电池}都会发热——
这正是\omterm{prop:g8:electric-current-effects:three}{热效应}的体现。
\end{solution}

\begin{solution}{exo:g8:electric-current-effects:6}
断开\omterm{ex:g3:first-electric-circuit:switch}{开关}：电流消失，磁性消失，汽车被放下。
若是永久\omterm{def:g1:magnets:magnet}{磁铁}，磁性强到这种程度，将永远吸住汽车——
要松开它，就得动用机器把它硬拆下来。
\end{solution}

\begin{solution}{exo:g8:electric-current-effects:7}
放电时：储存的化学能转化为电能，驱动电流流动，
最终转化为光能、机械能和内能。充电时，
电能被迫反向流过\omterm{def:g3:first-electric-circuit:battery}{电池}，重建
化学能储备。
\end{solution}

\begin{solution}{exo:g8:electric-current-effects:8}
因为电流本身看不见，只能通过它的效应来察觉。
发热、磁针偏转、\omterm{def:g3:solids-liquids-gases:liquid}{液体}冒泡，正是电流仅有的标志——
而下一章的仪器会把磁针的偏转精确为可读的
数值。
\end{solution}

\begin{solution}{exo:g8:electric-current-effects:9}
按下按钮：电流接通，\omterm{def:g8:electric-current-effects:electromagnet}{电磁铁}吸住锤子，发出“叮”的一声——
但锤子摆动时会断开自身电源中的触点；磁性消失，弹簧把锤子弹回，
触点重新接通；电流恢复，再次吸住锤子——如此反复，
每秒可达多次。\omterm{def:g8:electric-current-effects:electromagnet}{电磁铁}无法持续吸住锤子，这是设计使然：吸合
总会切断它自己的电流来源——于是连续的颤响
\emph{即是}
铃声。
\end{solution}

\begin{solution}{exo:g8:electric-current-effects:10}
恒温器中脆弱的触点只承载线圈的小电流；线圈的磁性吸合
粗壮的触点，这些触点在另一个独立回路中承载加热器的大电流。
小电流控制大电流，两个电路彼此从不接触——
精巧的操作只需拨动\omterm{def:g3:levers-and-scales:lever}{杠杆}，
而无须直接承受负载。
\end{solution}

\begin{solution}{exo:g8:electric-current-effects:11}
方法：将\omterm{def:g4:magnets-and-compass:compass}{指南针}平放在墙面上，远离铁质家具，
缓慢扫过墙面；磁针随可疑电器的\omterm{ex:g3:first-electric-circuit:switch}{开关}而偏转之处，
即有通电导线经过。
偏转越明显，说明电流越近或越强。
需要注意：铁质水管、\omterm{def:g1:magnets:magnet}{磁铁}乃至地球本身，都会使磁针偏转，
所以只有偏转随\omterm{ex:g3:first-electric-circuit:switch}{开关}的通断而同步变化，
才能证明存在电流。
\end{solution}

\begin{solution}{exo:g8:electric-current-effects:12}
电流使线圈变成\omterm{def:g1:magnets:magnet}{磁铁}；线圈与框架的异名\omterm{def:g4:magnets-and-compass:poles}{磁极}相互\omterm{def:g1:magnets:magnet}{吸引}，
同名\omterm{def:g4:magnets-and-compass:poles}{磁极}相互排斥，线圈便转向对齐的位置。
若不改变电流方向，线圈转到对齐位置便会停下——
转过半圈后抖动一下便归于静止，这只是门闭合器，而不是\omterm{ex:g6:circuits-and-safety:motor}{电动机}。
在线圈恰好对齐时改变电流方向，线圈的两极随之对调：
原先的\omterm{def:g1:magnets:magnet}{吸引}变为排斥，对齐位置反而成了最不利的位置，
线圈必须不断追逐一个每转过半圈就躲开的目标——
这种永无止境的追逐，
也就是转动。
\end{solution}

\begin{solution}{pb:g8:electric-current-effects:1}
\textbf{1.} 三个演示实验\emph{\omterm{def:g4:series-circuits:series}{串联}}在同一个带\omterm{ex:g3:first-electric-circuit:switch}{开关}的回路中：
同一股电流依次流经加热的导线、小磁针旁的导线和电解液——
一个\omterm{ex:g3:first-electric-circuit:switch}{开关}，三种效应。
\textbf{2.} \omterm{def:g4:series-circuits:series}{串联}规律：\emph{同一}电流无论流过粗导线还是细导线，大小都相同。
不同的只是路径——
细导线阻碍更大，因而发热更明显：
电流相同，受热各异。
\textbf{3.} 对调\omterm{def:g3:first-electric-circuit:battery}{电池}的正负极，\emph{并}把导线移到\omterm{def:g4:magnets-and-compass:compass}{指南针}的另一侧
（或将\omterm{def:g4:magnets-and-compass:compass}{指南针}调转）：两次反转相互抵消，
磁针的偏转方向仍指向东。
\textbf{4.} 水中两种成分以二比一的比例存在——
从一个电极得到的\omterm{def:g3:solids-liquids-gases:gas}{气体}，
\omterm{def:g6:mass-and-volume:volume}{体积}是另一个的两倍。
\textbf{5.} 去掉会自行断开的触点（即锤子摆动时断开的那个）：
\omterm{def:g8:electric-current-effects:electromagnet}{电磁铁}便会把毡制的锤子持续吸住，使其贴着铃而不发声——
参观者只能听到一声轻柔的“咚”，
此外再无声响。
\textbf{6.} 在废料箱旁，操作员断开线圈的\omterm{ex:g3:first-electric-circuit:switch}{开关}：
磁性消失，铁块落下。铝块留在原处，因为
无论天然\omterm{def:g1:magnets:magnet}{磁铁}还是电流产生的\omterm{def:g1:magnets:magnet}{磁铁}，都只\omterm{def:g1:magnets:magnet}{吸引}铁这一类物质，
铝从来不受它们\omterm{def:g1:magnets:magnet}{吸引}。
\textbf{7.} 电流产生的\omterm{def:g1:magnets:magnet}{磁铁}，其\omterm{def:g4:magnets-and-compass:poles}{磁极}随电流方向而定：
改变其中一个，另一个也随之改变。
\textbf{8.} \omterm{def:g3:first-electric-circuit:battery}{电池}靠\emph{消耗}自身储存的化学能来驱动电镀；
要重新储存化学能，需要从外部输入\omterm{def:g5:everyday-energy:energy}{能量}。
若\omterm{def:g3:first-electric-circuit:battery}{电池}能靠自己所做的功来补充自身\omterm{def:g5:everyday-energy:energy}{能量}，
便又回到了永动机的老问题——\omterm{def:g5:everyday-energy:energy}{能量}的收支
无法自己平衡自己。
\textbf{9.} 抓取：\omterm{ex:g3:first-electric-circuit:switch}{开关}闭合，线圈磁化，夹住小玩具。
摆动：电流保持，抓力保持——\omterm{prop:g8:electric-current-effects:three}{磁效应}正在起作用。
松开：\omterm{ex:g3:first-electric-circuit:switch}{开关}在滑槽上方断开，磁性消失，
奖品落下。
\textbf{10.} \omterm{def:g8:electric-current-effects:electromagnet}{电磁铁}依赖电流而存在：\omterm{def:g3:first-electric-circuit:battery}{电池}电压下降，电流减弱，
抓力随之减弱——可以随时通断，代价便是对电流的依赖。
（永久\omterm{def:g1:magnets:magnet}{磁铁}不会减弱，
也永远不会松开。）
\textbf{11.} 电流相同时，匝数越多，\omterm{def:g1:magnets:magnet}{磁铁}越强：
抓力越牢。另一种调节方式是电流本身的强弱——
即\emph{电流强度}，将在下一章中\omterm{def:g6:measurement-in-science:measuring}{测量}。
\textbf{12.} 例如：“\omterm{prop:g8:electric-current-effects:three}{热效应}：\omterm{def:g7:short-circuits-safety:fuse}{保险丝}，会率先忠实地熔断。
\omterm{prop:g8:electric-current-effects:three}{磁效应}：抓娃娃机，听从指令抓取与释放。
\omterm{prop:g8:electric-current-effects:three}{化学效应}：电镀槽，用\omterm{def:g3:first-electric-circuit:battery}{电池}的耐心
把勺子镀上银。”
\end{solution}
